\appendix

\chapter{Diccionario de datos final}

Este anexo resume las 12 variables que componen el dataset final integrado y deja explícito el origen de cada una.

\begin{table}[htbp]
\centering
\footnotesize
\caption{Diccionario resumido del dataset final integrado.}
\begin{tabular}{|p{3.2cm}|p{1.8cm}|p{5.5cm}|p{3.5cm}|}
\hline
Variable & Tipo & Descripción & Origen \\
\hline
codigo\_comuna & string(5) & Llave maestra comunal. & dim\_comuna\_base / todas las fuentes \\
nombre\_comuna & string & Nombre oficial de la base maestra. & dim\_comuna\_base \\
poblacion & int & Población censada comunal. & Fuente D / Censo 2024 \\
anio\_poblacion & int & Año de referencia de población. & Constante Fase 6 \\
pobreza\_ingresos\_pct & float & Porcentaje de pobreza por ingresos. & Fuente C / Observatorio Social \\
anio\_pobreza & int & Año de referencia de pobreza. & Constante Fase 6 \\
areas\_verdes\_m2 & int & Superficie total de áreas verdes comunales. & Fuente A / SINIM \\
anio\_areas\_verdes & int & Año de referencia de áreas verdes. & Constante Fase 6 \\
ipp\_miles\_pesos & int & Ingresos propios permanentes en miles de pesos. & Fuente B / SINIM \\
anio\_ingresos & int & Año de referencia del IPP. & Constante Fase 6 \\
areas\_verdes\_m2\_hab & float & Áreas verdes por habitante. & Derivada Fase 6 \\
ipp\_pesos\_hab & float & IPP en pesos por habitante. & Derivada Fase 6 \\
\hline
\end{tabular}
\end{table}


\chapter{Perfilado extendido}

El detalle extendido del perfilado se respalda en \texttt{outputs/perfilado\_fuentes.xlsx}. La siguiente tabla resume la estructura observada en la etapa diagnóstica.

\begin{table}[htbp]
\centering
\footnotesize
\caption{Perfilado resumido de las cuatro fuentes según el workbook de Fase 2.}
\begin{tabular}{|c|p{3.3cm}|r|r|r|r|p{4.0cm}|}
\hline
ID & Fuente & Filas & Cols. & Dup. & Filas nulas & Formato y parser \\
\hline
A & SINIM - Areas Verdes & 52 & 4 & 0 & 0 & xls (SpreadsheetML/XML Excel 2003) / spreadsheetml\_xml\_2003 \\
B & SINIM - Capacidad Municipal & 52 & 3 & 0 & 0 & xls (SpreadsheetML/XML Excel 2003) / spreadsheetml\_xml\_2003 \\
C & Observatorio Social - Pobreza por Ingresos & 351 & 10 & 0 & 1 & xlsx / pandas\_read\_excel \\
D & Censo 2024 - Poblacion Comunal & 349 & 10 & 0 & 1 & xlsx / pandas\_read\_excel \\
\hline
\end{tabular}
\end{table}


\chapter{Homologación de comunas}

La homologación se resolvió sin pendientes manuales dentro del universo final. Se documentan a continuación el resumen por fuente y ejemplos de normalización aplicada.

\begin{table}[htbp]
\centering
\footnotesize
\caption{Resumen por fuente de la homologación dentro del universo final.}
\begin{tabular}{|c|r|r|r|r|}
\hline
Fuente & Exactas & Por normalización & Manual & Pendientes \\
\hline
A & 26 & 6 & 0 & 0 \\
B & 26 & 6 & 0 & 0 \\
C & 0 & 32 & 0 & 0 \\
D & 0 & 32 & 0 & 0 \\
\hline
\end{tabular}
\end{table}


\begin{table}[htbp]
\centering
\footnotesize
\caption{Ejemplos reales de homologación resuelta por normalización automática.}
\begin{tabular}{|c|c|p{3.0cm}|p{3.0cm}|p{4.5cm}|}
\hline
Fuente & Código & Nombre base & Nombre observado & Resolución \\
\hline
A & 13104 & CONCHALI & CONCHALÍ & coincidencia por normalizacion \\
A & 13106 & ESTACION CENTRAL & ESTACIÓN CENTRAL & coincidencia por normalizacion \\
A & 13119 & MAIPU & MAIPÚ & coincidencia por normalizacion \\
A & 13122 & PEÑALOLEN & PEÑALOLÉN & coincidencia por normalizacion \\
A & 13129 & SAN JOAQUIN & SAN JOAQUÍN & coincidencia por normalizacion \\
A & 13131 & SAN RAMON & SAN RAMÓN & coincidencia por normalizacion \\
B & 13104 & CONCHALI & CONCHALÍ & coincidencia por normalizacion \\
B & 13106 & ESTACION CENTRAL & ESTACIÓN CENTRAL & coincidencia por normalizacion \\
\hline
\end{tabular}
\end{table}


\chapter{Mapa lógico extendido}

Este anexo deja trazabilidad entre artefactos del ETL y sus destinos operativos.

\begin{table}[htbp]
\centering
\footnotesize
\caption{Mapa lógico extendido de artefactos y destinos del ETL.}
\begin{tabular}{|p{4.3cm}|p{2.5cm}|p{5.2cm}|p{2.4cm}|}
\hline
Artefacto & Tipo & Contenido relevante & Fase de uso \\
\hline
data/raw/metadata\_fuentes.csv & Contrato operativo de lectura & Define archivo, hoja, skiprows y observaciones por fuente & Fases 2 y 4 \\
data/raw/dim\_comuna\_base.csv & Universo maestro & Restringe cobertura a 32 comunas y aporta nombre oficial & Fases 3, 5 y 6 \\
data/staging/poblacion\_staging.csv & Staging D & codigo\_comuna, nombre\_comuna, poblacion & Fase 6 / merge 1 \\
data/staging/pobreza\_staging.csv & Staging C & codigo\_comuna, nombre\_comuna, pobreza\_ingresos\_pct & Fase 6 / merge 2 \\
data/staging/areas\_verdes\_staging.csv & Staging A & codigo\_comuna, nombre\_comuna, mmpqc\_2024, mmpzc\_2024, areas\_verdes\_m2 & Fase 6 / merge 3 \\
data/staging/ingresos\_staging.csv & Staging B & codigo\_comuna, nombre\_comuna, ipp\_miles\_pesos & Fase 6 / merge 4 \\
data/processed/desigualdad\_comunal\_final.csv & Dataset final & 12 columnas finales y una fila por comuna & Fases 7, 8 y 9 \\
db/lab1\_desigualdad.sqlite & Persistencia final & dim\_comuna, fact\_desigualdad\_comunal, metadata\_fuentes & Fase 8 \\
\hline
\end{tabular}
\end{table}


\chapter{Validación final del dataset}

La validación de Fase 7 confirma consistencia estructural, cobertura y rangos.

\begin{table}[htbp]
\centering
\scriptsize
\caption{Resumen de validaciones estructurales y de rango del dataset final.}
\begin{tabular}{|p{2.4cm}|p{5.0cm}|c|p{3.2cm}|p{3.4cm}|}
\hline
Check & Descripción & Resultado & Valor observado & Detalle \\
\hline
fila\_por\_comuna & Existe una sola fila por comuna. & OK & filas=32; comunas\_unicas=32 & Sin multiplicacion de filas por comuna. \\
codigo\_comuna\_unico & `codigo\_comuna` es unico. & OK & 0 & Sin duplicados por llave. \\
cantidad\_comunas & El dataset final contiene exactamente 32 comunas. & OK & 32 & Cobertura total observada. \\
cobertura\_dim\_base & La cobertura comunal coincide con `dim\_comuna\_base.csv`. & OK & faltantes=[]; extras=[] & Cobertura exacta contra la base maestra. \\
convertibilidad\_numerica & Las columnas numericas son convertibles a tipo numerico razonable. & OK & poblacion=0; anio\_poblacion=0; pobreza\_ingresos\_pct=0; anio\_pobreza=0; areas\_verdes\_m2=0; anio\_areas\_verdes=0; ipp\_miles\_pesos=0; anio\_ingresos=0; areas\_verdes\_m2\_hab=0; ipp\_pesos\_hab=0 & Todas las columnas numericas son convertibles. \\
poblacion\_positiva & `poblacion` es estrictamente mayor que 0. & OK & 0 & Toda la poblacion es positiva. \\
pobreza\_rango & `pobreza\_ingresos\_pct` queda entre 0 y 100 cuando hay dato. & OK & 0 & La variable de pobreza esta dentro de rango. \\
areas\_verdes\_no\_negativas & `areas\_verdes\_m2` es no negativa cuando hay dato. & OK & 0 & Sin areas verdes negativas. \\
ipp\_no\_negativo & `ipp\_miles\_pesos` es no negativo cuando hay dato. & OK & 0 & Sin IPP negativos. \\
areas\_verdes\_m2\_hab\_sin\_inf & `areas\_verdes\_m2\_hab` no contiene infinitos. & OK & 0 & Sin infinitos en areas verdes por habitante. \\
ipp\_pesos\_hab\_sin\_inf & `ipp\_pesos\_hab` no contiene infinitos. & OK & 0 & Sin infinitos en IPP por habitante. \\
apto\_para\_sqlite & El dataset final es apto para pasar a Fase 8 / carga a SQLite. & OK & APTO & El dataset final cumple los criterios estructurales y de integridad requeridos. \\
\hline
\end{tabular}
\end{table}


\chapter{Carga SQLite}

Se resumen los conteos observados y las consultas de prueba utilizadas tras la carga final.

\begin{table}[htbp]
\centering
\footnotesize
\caption{Conteo de filas por tabla en la base SQLite final.}
\begin{tabular}{|p{4.2cm}|r|p{6.0cm}|}
\hline
Tabla & Filas & Rol \\
\hline
dim\_comuna & 32 & Dimensión comunal \\
fact\_desigualdad\_comunal & 32 & Tabla de hechos \\
metadata\_fuentes & 4 & Contrato operativo de fuentes \\
\hline
\end{tabular}
\end{table}


\begin{table}[htbp]
\centering
\scriptsize
\caption{Consultas de prueba documentadas tras la carga en SQLite.}
\begin{tabular}{|p{3.2cm}|p{0.9cm}|p{3.0cm}|p{1.2cm}|p{2.8cm}|p{2.3cm}|p{1.8cm}|}
\hline
Consulta & Ord. & Tabla & Código & Comuna & Valor & Detalle \\
\hline
conteo\_filas\_por\_tabla & - & dim\_comuna & - & - & 32.0 & conteo de filas \\
conteo\_filas\_por\_tabla & - & fact\_desigualdad\_comunal & - & - & 32.0 & conteo de filas \\
conteo\_filas\_por\_tabla & - & metadata\_fuentes & - & - & 4.0 & conteo de filas \\
top\_5\_menor\_areas\_verdes\_m2\_hab & 1.0 & fact\_desigualdad\_comunal & 13104 & CONCHALI & 0.603247 & areas\_verdes\_m2\_hab \\
top\_5\_menor\_areas\_verdes\_m2\_hab & 2.0 & fact\_desigualdad\_comunal & 13102 & CERRILLOS & 0.878035 & areas\_verdes\_m2\_hab \\
top\_5\_menor\_areas\_verdes\_m2\_hab & 3.0 & fact\_desigualdad\_comunal & 13108 & INDEPENDENCIA & 0.878941 & areas\_verdes\_m2\_hab \\
top\_5\_menor\_areas\_verdes\_m2\_hab & 4.0 & fact\_desigualdad\_comunal & 13109 & LA CISTERNA & 0.904776 & areas\_verdes\_m2\_hab \\
top\_5\_menor\_areas\_verdes\_m2\_hab & 5.0 & fact\_desigualdad\_comunal & 13130 & SAN MIGUEL & 1.154347 & areas\_verdes\_m2\_hab \\
top\_5\_mayor\_pobreza\_ingresos\_pct & 1.0 & fact\_desigualdad\_comunal & 13112 & LA PINTANA & 9.2938 & pobreza\_ingresos\_pct \\
top\_5\_mayor\_pobreza\_ingresos\_pct & 2.0 & fact\_desigualdad\_comunal & 13131 & SAN RAMON & 6.8821 & pobreza\_ingresos\_pct \\
top\_5\_mayor\_pobreza\_ingresos\_pct & 3.0 & fact\_desigualdad\_comunal & 13116 & LO ESPEJO & 6.7817 & pobreza\_ingresos\_pct \\
top\_5\_mayor\_pobreza\_ingresos\_pct & 4.0 & fact\_desigualdad\_comunal & 13105 & EL BOSQUE & 6.1982 & pobreza\_ingresos\_pct \\
top\_5\_mayor\_pobreza\_ingresos\_pct & 5.0 & fact\_desigualdad\_comunal & 13104 & CONCHALI & 6.0799 & pobreza\_ingresos\_pct \\
top\_5\_menor\_ipp\_pesos\_hab & 1.0 & fact\_desigualdad\_comunal & 13103 & CERRO NAVIA & 27701.980354 & ipp\_pesos\_hab \\
top\_5\_menor\_ipp\_pesos\_hab & 2.0 & fact\_desigualdad\_comunal & 13112 & LA PINTANA & 38281.739358 & ipp\_pesos\_hab \\
top\_5\_menor\_ipp\_pesos\_hab & 3.0 & fact\_desigualdad\_comunal & 13117 & LO PRADO & 39312.805346 & ipp\_pesos\_hab \\
top\_5\_menor\_ipp\_pesos\_hab & 4.0 & fact\_desigualdad\_comunal & 13105 & EL BOSQUE & 45941.619379 & ipp\_pesos\_hab \\
top\_5\_menor\_ipp\_pesos\_hab & 5.0 & fact\_desigualdad\_comunal & 13111 & LA GRANJA & 47063.353627 & ipp\_pesos\_hab \\
\hline
\end{tabular}
\end{table}


\chapter{Revisión breve de Fase 9: outliers y comunicabilidad}

La revisión corta de Fase 9 respalda tres advertencias: el índice de rezago es auxiliar, Quilicura debe tratarse como outlier plausible en áreas verdes por habitante y el scatter requiere lectura cuidadosa por su escala. Para el informe final conviene mantener estas advertencias junto a las figuras analíticas.

\begin{table}[htbp]
\centering
\footnotesize
\caption{Resumen de hallazgos descriptivos reutilizables desde los outputs de Fase 9.}
\begin{tabular}{|p{4.3cm}|p{7.0cm}|p{3.0cm}|}
\hline
Hallazgo & Comunas observadas & Output fuente \\
\hline
Menor áreas verdes por habitante & CONCHALI, CERRILLOS, INDEPENDENCIA, LA CISTERNA, SAN MIGUEL & outputs/ranking\_areas\_verdes.csv \\
Mayor pobreza por ingresos & LA PINTANA, SAN RAMON, LO ESPEJO, EL BOSQUE, CONCHALI & outputs/ranking\_pobreza.csv \\
Menor IPP por habitante & CERRO NAVIA, LA PINTANA, LO PRADO, EL BOSQUE, LA GRANJA & outputs/ranking\_ipp.csv \\
Peor quintil simultáneo de pobreza y áreas verdes & CONCHALI & outputs/tablas\_hallazgos\_fase9.csv \\
Peor quintil simultáneo de pobreza e IPP & LA PINTANA, SAN RAMON, LO ESPEJO, EL BOSQUE, LA GRANJA, CERRO NAVIA & outputs/tablas\_hallazgos\_fase9.csv \\
Rezago crítico en al menos dos dimensiones & LO ESPEJO, EL BOSQUE, CONCHALI, LA PINTANA, SAN RAMON, LA GRANJA, CERRO NAVIA & outputs/indice\_rezago\_territorial.csv \\
\hline
\end{tabular}
\end{table}

